% GENERATED FILE. Edit canonical lessons, then run npm run generate:books.
% canonical-source-hash: fnv1a64:f942edb3bcfc054a
% canonical-lessons: HI-C374-atkal, HI-C374-anuman, HI-C374-yaqin, HI-C374-amuman, HI-C374-buujhna, HI-R374-review-367-374

\chapter{Perhaps, Probably}
\label{ch:hi-perhaps-probably}
\hlchaptermodality{374}

\begin{chapteropening}
\textbf{By the end of this chapter:} \emph{I can guess what will happen if things go a certain way.}

Complete the last lesson of chapter 374: I can guess what will happen if things go a certain way.
\end{chapteropening}

\section[\texorpdfstring{\hi{अटकल} (aṭkal)}{अटकल (aṭkal)}]{\hi{अटकल} (aṭkal) — a guess, speculation}

\label{lesson:HI-C374-atkal}



\begin{quote}
Before the new one: say the Hindi for to miss, then the Hindi for useful.
\end{quote}

\subsection*{You'll want to know: \hi{अटकल}}
\textbf{\hi{अटकल}} — \emph{aṭkal} — \textquotedblleft{}a guess, speculation\textquotedblright{}.

\begin{quote}
\hi{लोगों} \hi{की} \hi{अटकल} \hi{है} \hi{कि} \hi{अगर} \hi{वह} \hi{चुनाव} \hi{जीत} \hi{गया} \hi{तो} \hi{मंत्री} \hi{बनेगा।}
\end{quote}

\begin{quote}
\emph{logoṁ kī aṭkal hai ki agar vah cunāv jīt gayā to mantrī banegā.}
\end{quote}

\begin{quote}
People speculate that if he wins the election, he'll become a minister.
\end{quote}

Feminine.

\begin{cousinweb}
\textbf{\hi{अटकल}} (\emph{aṭkal}) is inherited from \emph{*aṭṭakkalā}, \textquotedblleft{}a guess\textquotedblright{}, a form CDIAL reconstructs from its descendants (CDIAL 183).
\end{cousinweb}

\begin{grammarlens}[title={What you've built}]
One more word for saying what happens if.
\end{grammarlens}

\subsection*{Your turn}
\begin{itemize}
\raggedright
  \item \emph{Say it:} \emph{aṭkal}
  \item \emph{Say it:} \emph{aṭkal}, once more
  \item \emph{Say it:} \emph{logoṁ kī aṭkal hai ki agar vah cunāv jīt gayā to mantrī banegā.} — People speculate that if he wins the election, he'll become a minister.
  \item \emph{Recall:} say the Hindi for to miss, then the Hindi for useful, then say \emph{aṭkal} again
\end{itemize}

\begin{tcolorbox}[breakable,colback=teal!4,colframe=teal!35!black,title={Before you move on}]
What does \hi{अटकल} mean? (\textquotedblleft{}A guess, speculation\textquotedblright{}.) Say it once more.
\end{tcolorbox}
\section[\texorpdfstring{\hi{अनुमान} (anumān)}{अनुमान (anumān)}]{\hi{अनुमान} (anumān) — an estimate, a forecast}

\label{lesson:HI-C374-anuman}



\begin{quote}
Before the new one: say the Hindi for useful, then the Hindi for speculation.
\end{quote}

\subsection*{You'll want to know: \hi{अनुमान}}
\textbf{\hi{अनुमान}} — \emph{anumān} — \textquotedblleft{}an estimate, a forecast\textquotedblright{}.

\begin{quote}
\hi{अनुमान} \hi{के} \hi{अनुसार} \hi{अगर} \hi{ऐसी} \hi{गर्मी} \hi{रहेगी} \hi{तो} \hi{पानी} \hi{कम} \hi{पड़ेगा।}
\end{quote}

\begin{quote}
\emph{anumān ke anusār agar aisī garmī rahegī to pānī kam paṛegā.}
\end{quote}

\begin{quote}
According to the forecast, if the heat goes on like this, water will run short.
\end{quote}

Masculine.

\begin{cousinweb}
\textbf{\hi{अनुमान}} (\emph{anumān}) is a tatsama, borrowed from Sanskrit \emph{anumāna}, \textquotedblleft{}inferring, a conclusion\textquotedblright{} (Monier-Williams).
\end{cousinweb}

\begin{grammarlens}[title={What you've built}]
One more word for saying what happens if.
\end{grammarlens}

\subsection*{Your turn}
\begin{itemize}
\raggedright
  \item \emph{Say it:} \emph{anumān}
  \item \emph{Say it:} \emph{anumān}, once more
  \item \emph{Say it:} \emph{anumān ke anusār agar aisī garmī rahegī to pānī kam paṛegā.} — According to the forecast, if the heat goes on like this, water will run short.
  \item \emph{Recall:} say the Hindi for useful, then the Hindi for speculation, then say \emph{anumān} again
\end{itemize}

\begin{tcolorbox}[breakable,colback=teal!4,colframe=teal!35!black,title={Before you move on}]
What does \hi{अनुमान} mean? (\textquotedblleft{}An estimate, a forecast\textquotedblright{}.) Say it once more.
\end{tcolorbox}
\section[\texorpdfstring{\hi{यक़ीन} (yaqīn)}{यक़ीन (yaqīn)}]{\hi{यक़ीन} (yaqīn) — certainty, belief}

\label{lesson:HI-C374-yaqin}



\begin{quote}
Before the new one: say the Hindi for speculation, then the Hindi for a forecast.
\end{quote}

\subsection*{You'll want to know: \hi{यक़ीन}}
\textbf{\hi{यक़ीन}} — \emph{yaqīn} — \textquotedblleft{}certainty, belief\textquotedblright{}.

\begin{quote}
\hi{मुझे} \hi{यक़ीन} \hi{है} \hi{कि} \hi{अगर} \hi{तुम} \hi{रोज़} \hi{अभ्यास} \hi{करोगे} \hi{तो} \hi{ज़रूर} \hi{जीतोगे।}
\end{quote}

\begin{quote}
\emph{mujhe yaqīn hai ki agar tum roz abhyās karoge to zarūr jītoge.}
\end{quote}

\begin{quote}
I'm sure that if you practise every day, you'll win.
\end{quote}

Masculine.

\begin{cousinweb}
\textbf{\hi{यक़ीन}} (\emph{yaqīn}) is usually traced to Arabic \emph{yaqīn}, \textquotedblleft{}certainty\textquotedblright{}, which reached Hindi through Persian.
\end{cousinweb}

\begin{grammarlens}[title={What you've built}]
One more word for saying what happens if.
\end{grammarlens}

\subsection*{Your turn}
\begin{itemize}
\raggedright
  \item \emph{Say it:} \emph{yaqīn}
  \item \emph{Say it:} \emph{yaqīn}, once more
  \item \emph{Say it:} \emph{mujhe yaqīn hai ki agar tum roz abhyās karoge to zarūr jītoge.} — I'm sure that if you practise every day, you'll win.
  \item \emph{Recall:} say the Hindi for speculation, then the Hindi for a forecast, then say \emph{yaqīn} again
\end{itemize}

\begin{tcolorbox}[breakable,colback=teal!4,colframe=teal!35!black,title={Before you move on}]
What does \hi{यक़ीन} mean? (\textquotedblleft{}Certainty, belief\textquotedblright{}.) Say it once more.
\end{tcolorbox}
\section[\texorpdfstring{\hi{अमूमन} (amūman)}{अमूमन (amūman)}]{\hi{अमूमन} (amūman) — usually, generally}

\label{lesson:HI-C374-amuman}



\begin{quote}
Before the new one: say the Hindi for a forecast, then the Hindi for certainty.
\end{quote}

\subsection*{You'll want to know: \hi{अमूमन}}
\textbf{\hi{अमूमन}} — \emph{amūman} — \textquotedblleft{}usually, generally\textquotedblright{}.

\begin{quote}
\hi{अगर} \hi{आसमान} \hi{साफ़} \hi{होता} \hi{है} \hi{तो} \hi{अमूमन} \hi{रात} \hi{को} \hi{ठंड} \hi{होती} \hi{है।}
\end{quote}

\begin{quote}
\emph{agar āsmān sāf hotā hai to amūman rāt ko ṭhaṇḍ hotī hai.}
\end{quote}

\begin{quote}
If the sky is clear, it is usually cold at night.
\end{quote}

\begin{cousinweb}
\textbf{\hi{अमूमन}} (\emph{amūman}) is usually traced to Arabic \emph{ʿumūman}, \textquotedblleft{}generally\textquotedblright{}, which reached Hindi through Persian.
\end{cousinweb}

\begin{grammarlens}[title={What you've built}]
One more word for saying what happens if.
\end{grammarlens}

\subsection*{Your turn}
\begin{itemize}
\raggedright
  \item \emph{Say it:} \emph{amūman}
  \item \emph{Say it:} \emph{amūman}, once more
  \item \emph{Say it:} \emph{agar āsmān sāf hotā hai to amūman rāt ko ṭhaṇḍ hotī hai.} — If the sky is clear, it is usually cold at night.
  \item \emph{Recall:} say the Hindi for a forecast, then the Hindi for certainty, then say \emph{amūman} again
\end{itemize}

\begin{tcolorbox}[breakable,colback=teal!4,colframe=teal!35!black,title={Before you move on}]
What does \hi{अमूमन} mean? (\textquotedblleft{}Usually, generally\textquotedblright{}.) Say it once more.
\end{tcolorbox}
\section[\texorpdfstring{\hi{बूझना} (būjhnā)}{बूझना (būjhnā)}]{\hi{बूझना} (būjhnā) — to make out, to guess (a riddle)}

\label{lesson:HI-C374-buujhna}



\begin{quote}
Before the new one: say the Hindi for certainty, then the Hindi for usually.
\end{quote}

\subsection*{You'll want to know: \hi{बूझना}}
\textbf{\hi{बूझना}} — \emph{būjhnā} — \textquotedblleft{}to make out, to guess (a riddle)\textquotedblright{}.

\begin{quote}
\hi{अगर} \hi{तुमने} \hi{मेरी} \hi{पहेली} \hi{बूझ} \hi{ली} \hi{तो} \hi{मैं} \hi{हार} \hi{मान} \hi{लूँगा।}
\end{quote}

\begin{quote}
\emph{agar tumne merī pahelī būjh lī to maiṁ hār mān lūṅgā.}
\end{quote}

\begin{quote}
If you guess my riddle, I'll admit defeat.
\end{quote}

Past, with \hi{ने} (he or she …): \hi{उसने} \hi{पहेली} \hi{बूझ} \hi{ली}, \emph{usne pahelī būjh lī}.

\begin{cousinweb}
\textbf{\hi{बूझना}} (\emph{būjhnā}) is a tadbhava, inherited from Sanskrit \emph{budhyate}, \textquotedblleft{}wakes; understands\textquotedblright{} (CDIAL 9279).
\end{cousinweb}

\begin{grammarlens}[title={What you've built}]
One more word for saying what happens if.
\end{grammarlens}

\subsection*{Your turn}
\begin{itemize}
\raggedright
  \item \emph{Say it:} \emph{būjhnā}
  \item \emph{Say it:} \emph{būjhnā}, once more
  \item \emph{Say it:} \emph{agar tumne merī pahelī būjh lī to maiṁ hār mān lūṅgā.} — If you guess my riddle, I'll admit defeat.
  \item \emph{Recall:} say the Hindi for certainty, then the Hindi for usually, then say \emph{būjhnā} again
\end{itemize}

\begin{tcolorbox}[breakable,colback=teal!4,colframe=teal!35!black,title={Before you move on}]
What does \hi{बूझना} mean? (\textquotedblleft{}To make out, to guess (a riddle)\textquotedblright{}.) Say it once more.
\end{tcolorbox}
\section[(dialogue)]{Looking back: from Praise and Blame to Perhaps, Probably}

\label{lesson:HI-R374-review-367-374}



\begin{quote}
Say the Hindi for to commend and praise.
\end{quote}

\subsection*{Your turn}
Say these aloud:

\begin{itemize}
\raggedright
  \item to commend, praise, fault, criticism, to endure
  \item to get better, development, advancement, discipline, to polish
  \item corruption, a bribe (everyday), an institution, an employee, a ban
  \item an assembly, a village head, a district, a political party, to fix (a date)
  \item although, despite, instead of, apart from, to dig one's heels in
  \item mercy, compassion, loyal, a sin, to hold on to
  \item possible, possible (formal), a likelihood, to miss, useful
  \item speculation, a forecast, certainty, usually, to make out
\end{itemize}

\begin{tcolorbox}[breakable,colback=teal!4,colframe=teal!35!black,title={Before you move on}]
To commend? (\textbf{\hi{सराहना}}, \emph{sarāhnā}.) Praise? (\textbf{\hi{तारीफ़}}, \emph{tārīf}.) Fault? (\textbf{\hi{क़सूर}}, \emph{qasūr}.) Criticism? (\textbf{\hi{आलोचना}}, \emph{ālocnā}.) To endure? (\textbf{\hi{झेलना}}, \emph{jhelnā}.) To get better? (\textbf{\hi{सुधरना}}, \emph{sudharnā}.) Development? (\textbf{\hi{विकास}}, \emph{vikās}.) Advancement? (\textbf{\hi{उन्नति}}, \emph{unnati}.) Discipline? (\textbf{\hi{अनुशासन}}, \emph{anuśāsan}.) To polish? (\textbf{\hi{निखारना}}, \emph{nikhārnā}.) Corruption? (\textbf{\hi{भ्रष्टाचार}}, \emph{bhraṣṭācār}.) A bribe (everyday)? (\textbf{\hi{घूस}}, \emph{ghūs}.) An institution? (\textbf{\hi{संस्था}}, \emph{saṁsthā}.) An employee? (\textbf{\hi{कर्मचारी}}, \emph{karmacārī}.) A ban? (\textbf{\hi{पाबंदी}}, \emph{pābandī}.) An assembly? (\textbf{\hi{सभा}}, \emph{sabhā}.) A village head? (\textbf{\hi{प्रधान}}, \emph{pradhān}.) A district? (\textbf{\hi{ज़िला}}, \emph{zilā}.) A political party? (\textbf{\hi{पार्टी}}, \emph{pārṭī}.) To fix (a date)? (\textbf{\hi{तय} \hi{करना}}, \emph{tay karnā}.) Although? (\textbf{\hi{हालाँकि}}, \emph{hālā̃ki}.) Despite? (\textbf{\hi{बावजूद}}, \emph{bāvjūd}.) Instead of? (\textbf{\hi{बजाय}}, \emph{bajāy}.) Apart from? (\textbf{\hi{अलावा}}, \emph{alāvā}.) To dig one's heels in? (\textbf{\hi{अड़ना}}, \emph{aṛnā}.) Mercy? (\textbf{\hi{दया}}, \emph{dayā}.) Compassion? (\textbf{\hi{करुणा}}, \emph{karuṇā}.) Loyal? (\textbf{\hi{वफ़ादार}}, \emph{vafādār}.) A sin? (\textbf{\hi{पाप}}, \emph{pāp}.) To hold on to? (\textbf{\hi{थामना}}, \emph{thāmnā}.) Possible? (\textbf{\hi{मुमकिन}}, \emph{mumkin}.) Possible (formal)? (\textbf{\hi{संभव}}, \emph{sambhav}.) A likelihood? (\textbf{\hi{संभावना}}, \emph{sambhāvnā}.) To miss? (\textbf{\hi{चूकना}}, \emph{cūknā}.) Useful? (\textbf{\hi{उपयोगी}}, \emph{upayogī}.) Speculation? (\textbf{\hi{अटकल}}, \emph{aṭkal}.) A forecast? (\textbf{\hi{अनुमान}}, \emph{anumān}.) Certainty? (\textbf{\hi{यक़ीन}}, \emph{yaqīn}.) Usually? (\textbf{\hi{अमूमन}}, \emph{amūman}.) To make out? (\textbf{\hi{बूझना}}, \emph{būjhnā}.)
\end{tcolorbox}
